% Figura autónoma de paquetes: cargar tikz y arrows.meta en el preámbulo.
% Dimensiones naturales: aproximadamente 26 x 17 cm; usar página apaisada.
% Las operaciones se seleccionaron de los encabezados del firmware.
\begin{tikzpicture}[
  x=1cm,y=1cm,
  font=\sffamily\fontsize{9}{11}\selectfont,
  call/.style={-{Latex[length=2.1mm,width=1.5mm]},draw=archBlue,line width=.75pt,
    rounded corners=2pt,preaction={draw=white,line width=3pt}},
  support/.style={call,draw=archGray,dashed,line width=.6pt},
  cb/.style={call,draw=archOrange,densely dotted,line width=1pt},
  tag/.style={font=\sffamily\fontsize{8}{9}\selectfont,fill=white,
    text=archInk,inner sep=2pt,align=center},
  op/.style={anchor=north west,inner sep=0pt,align=left,
    font=\ttfamily\fontsize{8.6}{11.5}\selectfont,text=archInk}
]
\definecolor{archInk}{HTML}{203247}
\definecolor{archBlue}{HTML}{365E88}
\definecolor{archTeal}{HTML}{287A79}
\definecolor{archPurple}{HTML}{766094}
\definecolor{archOrange}{HTML}{B87523}
\definecolor{archGray}{HTML}{7C8793}
\definecolor{archPaper}{HTML}{F5F7FA}

% Caja: identificador, x, y, altura, color, título, responsabilidad, operaciones.
\newcommand{\archmodule}[8]{%
  \coordinate (#1-nw) at (#2,#3);
  \coordinate (#1-ne) at (#2+7.4,#3);
  \coordinate (#1-sw) at (#2,#3-#4);
  \coordinate (#1-se) at (#2+7.4,#3-#4);
  \path[draw=#5!55,line width=.65pt,fill=white,rounded corners=3pt]
    (#1-nw) rectangle (#1-se);
  \path[fill=#5!9] (#2+.015,#3-.015) rectangle (#2+7.385,#3-1.03);
  \node[anchor=north west,inner sep=0pt,text=#5,
    font=\sffamily\fontsize{7.5}{9}\selectfont]
    at (#2+.22,#3-.14) {MÓDULO C};
  \node[anchor=north west,inner sep=0pt,text=archInk,
    font=\ttfamily\bfseries\fontsize{12}{14}\selectfont]
    at (#2+.22,#3-.39) {#6};
  \node[anchor=north west,inner sep=0pt,text=archInk,
    font=\sffamily\fontsize{8.2}{10}\selectfont]
    at (#2+.22,#3-.77) {#7};
  \draw[#5!35,line width=.5pt] (#2,#3-1.03)--(#2+7.4,#3-1.03);
  \node[op,text width=7cm] at (#2+.22,#3-1.23) {#8};
}

% Fila superior: coordinación e interfaces de aplicación.
\archmodule{parser}{0}{0}{3.25}{archPurple}{gcode\_parser}
  {Interpretación de una línea de G-code}
  {+ gcode\_parse\_line(...)\\[5pt]
   \normalfont\sffamily Salida: estructura \texttt{GcodeCommand}.\\
   \normalfont\sffamily La conversión numérica es interna.}

\archmodule{main}{9}{0}{3.8}{archBlue}{main.c}
  {UART, comandos, segmentos pendientes y alarmas}
  {+ main()\\
   - process\_gcode\_line(...)\\
   - flush\_pending\_move(...)\\
   - service\_realtime()}

\archmodule{ena}{18}{0}{3.25}{archTeal}{motor\_enable}
  {Habilitación común de los drivers}
  {+ motor\_enable\_init()\\
   + motors\_enable()\\
   + motors\_disable()\\
   + motors\_are\_enabled()}

% Fila central: referenciación y movimiento.
\archmodule{homing}{0}{-5.5}{3.5}{archOrange}{homing}
  {Referenciación Z, X e Y; lectura de límites}
  {+ homing\_init()\\
   + homing\_run()\\
   + homing\_is\_active()\\[5pt]
   \normalfont\sffamily Búsqueda directa: pulsos GPIO y DWT.}

\archmodule{step}{9}{-5.5}{4.9}{archTeal}{stepper}
  {Movimiento CoreXY e interrupción de TIM2}
  {+ stepper\_init(...)\\
   + stepper\_move\_to(...)\\
   + stepper\_get\_position(...)\\
   + stepper\_set\_position(...)\\
   + stepper\_emergency\_stop()\\
   + stepper\_timer\_isr()\\
   + stepper\_set\_wait\_callback(...) }

\archmodule{z}{18}{-5.5}{4.3}{archTeal}{z\_axis}
  {Movimiento Z e interrupción de TIM4}
  {+ z\_axis\_init(...)\\
   + z\_axis\_move\_to(...)\\
   + z\_axis\_get\_position(...)\\
   + z\_axis\_set\_position(...)\\
   + z\_axis\_emergency\_stop()\\
   + z\_axis\_timer\_isr()}

% Fila inferior: servicio compartido, láser y leyenda.
\archmodule{wd}{0}{-12}{2.65}{archGray}{watchdog}
  {Supervisión mediante IWDG}
  {+ watchdog\_refresh()\\
   + watchdog\_caused\_last\_reset()}

\archmodule{laser}{9}{-12}{4.3}{archTeal}{laser}
  {Accionamiento PWM mediante TIM3}
  {+ laser\_init(...)\\
   + laser\_enable() / laser\_disable()\\
   + laser\_set\_power(...)\\
   + laser\_set\_dynamic(...)\\
   + laser\_begin\_segment(...)\\
   + laser\_set\_speed\_ratio(...)}

% Dependencias directas de main.c.
\draw[call] (9,-1.8)--node[tag,above]{analiza}(7.4,-1.8);
\draw[call] (16.4,-1.8)--node[tag,above]{habilita}(18,-1.8);
\draw[call] (10,-3.8)--(10,-4.4)--(3.7,-4.4)--(3.7,-5.5);
\draw[call] (12.7,-3.8)--node[tag,right]{ejecuta}(12.7,-5.5);
\draw[call] (15.3,-3.8)--(15.3,-4.4)--(21.7,-4.4)--(21.7,-5.5);
\draw[call] (16.4,-2.7)--(17.35,-2.7)--(17.35,-11.45)--(15.4,-11.45)--(15.4,-12);
\node[tag,rotate=90] at (17.35,-8.1) {modo, potencia y segmento};

% Colaboraciones entre módulos de movimiento y referenciación.
\draw[call] (7.4,-7.2)--node[tag,above,text width=1.35cm]{mueve /\\fija cero}(9,-7.2);
\draw[call] (7.4,-8.3)--(7.85,-8.3)--(7.85,-10.95)--(25.95,-10.95)--(25.95,-8.5)--(25.4,-8.5);
\node[tag] at (20.6,-10.95) {referenciación Z};
\draw[call] (12.7,-10.4)--node[tag,left]{relación de\\velocidades}(12.7,-12);

% La devolución del control se produce por una función registrada, no por include.
\draw[cb] (16.4,-6.25)--(16.85,-6.25)--(16.85,-3.25)--(16.4,-3.25);
\node[tag,rotate=90,text=archOrange] at (16.85,-4.9) {callback};

% Supervisión: dependencias reales hacia watchdog.
\draw[support] (9,-3.25)--(8.4,-3.25)--(8.4,-3.65)--(-.5,-3.65)--(-.5,-11.55)--(6.3,-11.55)--(6.3,-12);
\draw[support] (2.2,-9)--node[tag,left]{refresca}(2.2,-12);
\draw[support] (20.4,-9.8)--(20.4,-10.55)--(5,-10.55)--(5,-12);

% Leyenda y límites de la representación.
\path[fill=archPaper,rounded corners=3pt] (18,-12) rectangle (25.4,-16.3);
\node[anchor=north west,text=archInk,font=\sffamily\bfseries\small]
  at (18.2,-12.15) {Cómo leer la figura};
\draw[call] (18.25,-12.83)--(19.05,-12.83);
\node[anchor=west,font=\sffamily\fontsize{8.3}{10}\selectfont,text=archInk]
  at (19.15,-12.83) {Llamada a una interfaz};
\draw[support] (18.25,-13.3)--(19.05,-13.3);
\node[anchor=west,font=\sffamily\fontsize{8.3}{10}\selectfont,text=archInk]
  at (19.15,-13.3) {Servicio de supervisión};
\draw[cb] (18.25,-13.77)--(19.05,-13.77);
\node[anchor=west,font=\sffamily\fontsize{8.3}{10}\selectfont,text=archInk]
  at (19.15,-13.77) {Invocación de callback};
\node[anchor=north west,text width=6.9cm,align=left,text=archInk,
  font=\sffamily\fontsize{8.2}{10.5}\selectfont] at (18.25,-14.15)
  {\texttt{+} Interfaz expuesta por el módulo.\\
   \texttt{-} Operación interna de \texttt{main.c}.\\[4pt]
   Se muestran operaciones seleccionadas y firmas abreviadas.
   La callback se ejecuta durante la espera del movimiento X--Y;
   \texttt{stepper} no depende de un encabezado de \texttt{main}.};
\end{tikzpicture}
